\documentclass[11pt]{book}
\usepackage[paperwidth=145mm,paperheight=200mm,top=18mm,bottom=20mm,inner=20mm,outer=16mm]{geometry}
\usepackage{brumfiel}
% figures13.tex -- Chapter 13 (Loci and Sets) figures.
%
% Constrained points are DERIVED, never placed by hand:
%   parallels    -> equal ratios along the two sides   ($(X)!r!(Y)$)
%   perp. feet   -> ($(X)!(P)!(Y)$)
%   crossings    -> (intersection of A--B and C--D)
%   bisectors    -> ratio from the angle-bisector theorem
%   midpoints    -> ($(X)!0.5!(Y)$)
% so the hypothesis stated in the book holds exactly in the drawing.
% Hypotheses are re-asserted numerically in tools/constraints/ch13.py.
%
% Macro names are chapter-prefixed: \FIGXIII<NUMBERWORD>.
%
% Line treatment follows the book's own convention: solid where the book
% prints solid (given), dashed where the book prints dashed (aux).  In this
% chapter the dashed role is doing double duty, exactly as the book does it:
%   - a LOCUS that is being asserted rather than given (Figs 13-1, 13-2)
%   - auxiliary/proof lines (Figs 13-5 ... 13-7, 13-11, 13-12, 13-14)
%   - the second position of a rolling wheel (Fig 13-9)
%   - the constructed outer triangle A'B'C' (Fig 13-16)
% The book shades no region anywhere in this chapter, so no fill patterns.

\usetikzlibrary{arrows.meta}

% Primed labels (Fig 13-16).  Same 1.3pt drop as ch09: keeps the prime in the
% same "word" as its letter for the collision audit.
\newcommand{\XIIIpr}{\raisebox{-1.3pt}{\ensuremath{{}^{\prime}}}}

% Chapter-local label clearance, as ch09: 3.4pt clips when a labelled point
% lies ON a drawn line.
\tikzset{
  vlab/.append style={outer sep=2.9pt},
  slab/.append style={outer sep=2.9pt},
}

% Long, widely spaced dashes.  Caleb, 2026-08-17: in Figs 13-11/13-12/13-16 the
% standard `aux' pattern is so fine that at print size the dashed cevians read
% as solid and the figure turns into "a quadrilateral with diagonals".  Measured
% off the print, the book's dashes in those three figures run about 2.6 mm on
% and 1.4 mm off -- roughly double ours.
\tikzset{auxlong/.style={aux, dash pattern=on 6.6pt off 3.8pt}}

% Small filled arrowheads, at a FIXED printed size so they do not grow with the
% stroke weight or shrink with the picture's `scale='.  The book's arrowheads
% are small; ours were blotting where several met at one vertex (Fig 13-16).
\tikzset{
  bkarrow/.style  = {-{Latex[length=3.4pt,width=2.8pt]}},
  bkarrows/.style = {{Latex[length=3.4pt,width=2.8pt]}-{Latex[length=3.4pt,width=2.8pt]}},
}

% The book marks the concurrency point in Figs 13-11, 13-12 and the centre in
% Fig 13-14 with a small saltire/plus rather than a filled dot.  In 13-11 and
% 13-12 that mark is a bold six-rayed asterisk, clearly heavier than the lines
% that cross there -- \XIIIstar reproduces it.
\newcommand{\XIIIcross}[1]{%
  \draw[fig,line width=0.5pt] ($(#1)+(-0.10,-0.10)$) -- ($(#1)+(0.10,0.10)$);
  \draw[fig,line width=0.5pt] ($(#1)+(-0.10,0.10)$) -- ($(#1)+(0.10,-0.10)$);}
\newcommand{\XIIIplus}[1]{%
  \draw[fig,line width=0.5pt] ($(#1)+(-0.11,0)$) -- ($(#1)+(0.11,0)$);
  \draw[fig,line width=0.5pt] ($(#1)+(0,-0.11)$) -- ($(#1)+(0,0.11)$);}
\newcommand{\XIIIstar}[1]{%
  \foreach \a in {15,75,135}
    \draw[fig,line width=0.75pt] ($(#1)+(\a:0.15)$) -- ($(#1)+(\a+180:0.15)$);}

% ---------------------------------------------------------------- Fig 13-1
% Line l solid; the locus (a pair of parallels at distance d) dashed, one on
% each side.  The d-tick is drawn from l up to the upper parallel.
\newcommand{\FIGXIIIONE}{\begin{bkfigure}[\linewidth]
\begin{tikzpicture}[scale=1.300]
  \draw[given] (0,0) -- (6.10,0);
  \node[vlab, right] at (6.10,0) {$l$};
  \draw[aux] (0.28,0.82) -- (5.82,0.82);
  \draw[aux] (0.28,-0.82) -- (5.82,-0.82);
  % the distance tick, with the book's short cap at the top
  \draw[given] (1.62,0) -- (1.62,0.82);
  \draw[given] (1.50,0.82) -- (1.74,0.82);
  \node[slab, right] at (1.66,0.42) {$d$};
\end{tikzpicture}
\figcap{13--1}\end{bkfigure}}

% ---------------------------------------------------------------- Fig 13-2
% Two given parallels solid; the locus (the midline) dashed.  P sits on the
% midline, and the transversal tick carries d above and d below.
\newcommand{\FIGXIIITWO}{\begin{bkfigure}[\linewidth]
\begin{tikzpicture}[scale=1.300]
  \coordinate (T) at (0,0.86);
  \coordinate (Bm) at (0,-0.86);
  \draw[given] (0,0.86) -- (6.00,0.86);
  \draw[given] (0,-0.86) -- (6.00,-0.86);
  \draw[aux] (0.22,0) -- (5.78,0);
  % transversal: P is the midpoint of the cut, so the two d's are equal
  \coordinate (Pt) at (3.46,0.86);
  \coordinate (Pb) at (3.46,-0.86);
  \coordinate (P)  at ($(Pt)!0.5!(Pb)$);
  \draw[given] (Pt) -- (Pb);
  \node[vlab, below left] at (P) {$P$};
  \node[slab, right] at (3.50,0.43) {$d$};
  \node[slab, right] at (3.50,-0.43) {$d$};
\end{tikzpicture}
\figcap{13--2}\end{bkfigure}}

% ---------------------------------------------------------------- Fig 13-3
% Ex 13-1 #2: a 3-unit segment, ticked into units, with the locus only
% partly sketched -- five dots at height 1 over the segment, and one dot at
% distance 1 measured radially off the right endpoint.  Reproduced as drawn.
\newcommand{\FIGXIIITHREE}{\begin{bkfigure}[\linewidth]
\begin{tikzpicture}[scale=1.196]
  \draw[given] (0,0) -- (3,0);
  \foreach \x in {0,1,2,3} \draw[given] (\x,0) -- (\x,0.17);
  % five dots, spacing measured from the print (0.19 unit); the arrow lands on
  % the FOURTH dot, with three dots to its left and one to its right, as drawn
  \foreach \x in {1.48,1.67,1.86,2.05,2.24} \fill (\x,1.0) circle (0.052);
  \draw[given,<->] (2.05,0.06) -- (2.05,0.94);
  \node[slab, left] at (2.01,0.50) {$1$};
  % radial unit off the right end: 72 degrees, length exactly 1
  \coordinate (R) at (3,0);
  \coordinate (Rd) at ($(R)+(72:1)$);
  \fill (Rd) circle (0.052);
  \draw[given,<->] ($(R)!0.07!(Rd)$) -- ($(R)!0.93!(Rd)$);
  \node[slab, right] at ($(R)!0.5!(Rd)$) {$1$};
\end{tikzpicture}
\figcap{13--3}\end{bkfigure}}

% ---------------------------------------------------------------- Fig 13-4
% Ex 13-1 #5: fixed base of 4 units, free vertex V.
\newcommand{\FIGXIIIFOUR}{\begin{bkfigure}[\linewidth]
\begin{tikzpicture}[scale=1.196]
  \coordinate (L) at (0,0);
  \coordinate (R) at (3.00,0);
  % apex measured off the print: 0.741 of the base along, 0.463 of the base up
  \coordinate (V) at (2.223,1.389);
  \draw[given] (L) -- (R) -- (V) -- cycle;
  \node[vlab, above] at (V) {$V$};
  \node[slab, below] at ($(L)!0.5!(R)$) {$4$};
\end{tikzpicture}
\figcap{13--4}\end{bkfigure}}

% ---------------------------------------------------------------- Fig 13-5
% Thm 13-1.  C is the MIDPOINT of AB; the perpendicular bisector is solid and
% runs through C to P; PA and PB are the dashed proof lines.
\newcommand{\FIGXIIIFIVE}{\begin{bkfigure}[\linewidth]
\begin{tikzpicture}[scale=1.300]
  \coordinate (A) at (0,0);
  \coordinate (B) at (3.40,0);
  \coordinate (C) at ($(A)!0.5!(B)$);
  \coordinate (P) at ($(C)+(0,1.56)$);
  \draw[given] (A) -- (B);
  % the bisector overshoots P by about 0.15 of AB in the print, and drops
  % about 0.13 of AB below the segment
  \draw[given] ($(C)+(0,-0.44)$) -- ($(C)+(0,2.06)$);
  \draw[aux] (A) -- (P);
  \draw[aux] (P) -- (B);
  \fill (C) circle (0.055);
  \node[vlab, below left] at (A) {$A$};
  \node[vlab, below right] at (C) {$C$};
  \node[vlab, below right] at (B) {$B$};
  \node[vlab, above right] at (P) {$P$};
\end{tikzpicture}
\figcap{13--5}\end{bkfigure}}

% ---------------------------------------------------------------- Fig 13-6
% Thm 13-2.  The sides of the angle are solid; the bisector is dashed and
% carries P.  A and B are the FEET of the perpendiculars from P, so they are
% built with the projection operator -- never by eye.
% Re-drawn 2026-08-17 against the print (Caleb's photo of p.229).  Four things
% changed: (a) the arrowhead that pointed INTO the vertex is gone -- the book
% has no arrowhead anywhere in this figure; (b) the dashed bisector now stops a
% little past P instead of running to the edge of the picture; (c) PA and PB are
% two straight segments meeting in a visible corner at P, with PA dropping
% VERTICALLY to A (the book lays the lower side of the angle horizontal, which
% is what makes that read); (d) the book's short cross-strokes at the two feet
% and at P are drawn.  Opening at O is 32 degrees, as in the print.
\newcommand{\FIGXIIISIX}{\begin{bkfigure}[\linewidth]
\begin{tikzpicture}[scale=1.300]
  \coordinate (O) at (0,0);
  \coordinate (Uend) at (32:3.56);   % upper side of the angle
  \coordinate (Lend) at (0:3.56);    % lower side, horizontal as in the print
  % |OP| = 2.85/cos 16 puts each foot at 0.80 of its drawn side, as the book does
  \coordinate (P) at (16:2.9648);    % on the bisector of the 32-degree angle
  \coordinate (B) at ($(O)!(P)!(Uend)$);
  \coordinate (A) at ($(O)!(P)!(Lend)$);
  \draw[given] (O) -- (Uend);
  \draw[given] (O) -- (Lend);
  \draw[aux] (O) -- ($(O)!1.25!(P)$);
  \draw[aux] (B) -- (P);
  \draw[aux] (P) -- (A);
  % the print's cross-strokes: one at each perpendicular foot, one at P
  \draw[tick] ($(A)!-0.13!(P)$) -- ($(A)!0.13!(P)$);
  \draw[tick] ($(B)!-0.13!(P)$) -- ($(B)!0.13!(P)$);
  \draw[tick] ($(P)+(106:0.13)$) -- ($(P)+(-74:0.13)$);
  \node[vlab, below left] at (O) {$O$};
  \node[vlab, above left] at (B) {$B$};
  \node[vlab] at ($(A)+(-84:0.32)$) {$A$};
  \node[vlab] at ($(P)+(62:0.34)$) {$P$};
\end{tikzpicture}
\figcap{13--6}\end{bkfigure}}

% ---------------------------------------------------------------- Fig 13-7
% Thm 13-3.  AB is a DIAMETER (O the midpoint); P on the circle; PA, PB dashed
% so that angle APB is the right angle in a semicircle.
\newcommand{\FIGXIIISEVEN}{\begin{bkfigure}[\linewidth]
\begin{tikzpicture}[scale=1.300]
  \def\rr{1.42}
  \coordinate (O) at (0,0);
  \coordinate (A) at (-\rr,0);
  \coordinate (B) at (\rr,0);
  \coordinate (P) at (118:\rr);
  \draw[given] (O) circle (\rr);
  \draw[given] (A) -- (B);
  \draw[aux] (P) -- (A);
  \draw[aux] (P) -- (B);
  \fill (O) circle (0.055);
  \node[vlab, left] at (A) {$A$};
  \node[vlab, right] at (B) {$B$};
  \node[vlab, below] at (O) {$O$};
  \node[vlab, above left] at (P) {$P$};
\end{tikzpicture}
\figcap{13--7}\end{bkfigure}}

% ---------------------------------------------------------------- Fig 13-8
% Ex 13-2 #4: A and B fixed on a dashed line; P below, on one side, with the
% solid segments PA, PB subtending the 50-degree angle of the exercise.
\newcommand{\FIGXIIIEIGHT}{\begin{bkfigure}[\linewidth]
\begin{tikzpicture}[scale=1.118]
  % AB rises at 29 degrees, as measured off the print (PDF 244 at 200 dpi).
  \coordinate (A) at (0.62,0.50);
  \coordinate (B) at (2.3255,1.4454);
  % P is placed on the 50-degree arc over AB, so that angle APB is EXACTLY the
  % 50 degrees the exercise states.  Measured against the print, the drawn P
  % already subtends 50.1 degrees, and it sits at 0.062 of AB along and 0.903
  % of AB below the line -- the position reproduced here.
  \coordinate (P) at (1.5799,-0.9823);
  \draw[aux] ($(A)!-0.68!(B)$) -- ($(B)!-1.00!(A)$);
  \draw[given] (A) -- (P);
  \draw[given] (B) -- (P);
  \fill (B) circle (0.055);
  \node[vlab, above left] at (A) {$A$};
  \node[vlab, above left] at (B) {$B$};
  \node[vlab, below] at (P) {$P$};
\end{tikzpicture}
\figcap{13--8}\end{bkfigure}}

% ---------------------------------------------------------------- Fig 13-9
% Ex 13-2 #6: the wheel in two positions -- the second dashed -- with the
% tracing point P on the rim.  Equal radii; the centres sit at the same
% height, which is what rolling on a (horizontal) line forces.
% Re-drawn 2026-08-17.  The two positions of the wheel are the SAME size and are
% offset SIDEWAYS, so they cross cleanly near the top and near the bottom; the
% old drawing let them all but coincide, which read as one circle nested inside
% another.  The short solid connector that the book draws -- the first piece of
% the traced path, from P's earlier position on the dashed rim to its later
% position on the solid rim, with a cross-tick at the earlier end -- was missing
% entirely, leaving the P label floating with nothing to name.
% Offsets measured off the print: centre separation 0.43 r (= 0.215 diameter),
% P at 241 degrees on the earlier wheel.
\newcommand{\FIGXIIININE}{\begin{bkfigure}[\linewidth]
\begin{tikzpicture}[scale=1.300]
  \def\rr{1.05}
  \def\dd{0.4515}                 % = 0.43\rr, how far the wheel has rolled
  \coordinate (Kp) at (0,0);      % dashed wheel, EARLIER position
  \coordinate (K)  at (\dd,0);    % solid wheel, LATER position; centres level
  % rolling without slipping turns the wheel through d/r = 0.43 rad = 24.637 deg
  \coordinate (P0) at ($(Kp)+(241:\rr)$);
  \coordinate (P)  at ($(K)+(216.363:\rr)$);
  \draw[aux] (Kp) circle (\rr);
  \draw[given] (K) circle (\rr);
  \draw[given] (P0) -- (P);
  \draw[tick] ($(P0)!0.33!90:(P)$) -- ($(P0)!-0.33!90:(P)$);
  \node[vlab] at ($(P)+(38:0.36)$) {$P$};
\end{tikzpicture}
\figcap{13--9}\end{bkfigure}}

% --------------------------------------------------------------- Fig 13-10
% Ex 13-2 #8: a circle of radius R/4 rolling INSIDE a circle of radius R.
% Internal tangency is forced: centre distance = R - R/4.
\newcommand{\FIGXIIITEN}{\begin{bkfigure}[\linewidth]
\begin{tikzpicture}[scale=1.300]
  \def\RR{1.40}
  \coordinate (O) at (0,0);
  \coordinate (Q) at (-32:0.75*\RR);   % = R - R/4 = 0.75 R
  \draw[given] (O) circle (\RR);
  \draw[given] (Q) circle (0.25*\RR);
  \fill (O) circle (0.05);
  \fill (Q) circle (0.05);
\end{tikzpicture}
\figcap{13--10}\end{bkfigure}}

% --------------------------------------------------------------- Fig 13-11
% Thm 13-4, INCENTER.  Each bisector foot is placed by the angle-bisector
% theorem (ratio of the two adjacent sides), so the three dashed cevians are
% genuinely the angle bisectors and genuinely concurrent.
\newcommand{\FIGXIIIELEVEN}{\begin{bkfigure}[\linewidth]
\begin{tikzpicture}[scale=1.300]
  \coordinate (A) at (0,0);
  \coordinate (B) at (3.50,2.70);
  \coordinate (C) at (3.30,0.10);
  % |BC| = 2.6077, |CA| = 3.3015, |AB| = 4.4204
  % foot on BC from A divides BC as AB : AC = 4.4204 : 3.3015
  \coordinate (Da) at ($(B)!0.57243!(C)$);
  % foot on CA from B divides AC as AB : BC = 4.4204 : 2.6077
  \coordinate (Eb) at ($(A)!0.62896!(C)$);
  % foot on AB from C divides AB as AC : BC = 3.3015 : 2.6077
  \coordinate (Fc) at ($(A)!0.55871!(B)$);
  \coordinate (O) at (intersection of A--Da and B--Eb);
  \draw[given] (A) -- (B) -- (C) -- cycle;
  % long dashes (see auxlong above): at the standard pattern these three read
  % as solid and the figure turns into a quadrilateral with its diagonals
  \draw[auxlong] (A) -- (Da);
  \draw[auxlong] (B) -- (Eb);
  \draw[auxlong] (C) -- (Fc);
  % the print marks the concurrency with a bold asterisk, not a bare crossing
  \XIIIstar{O}
  \node[vlab, below left] at (A) {$A$};
  \node[vlab, above] at (B) {$B$};
  \node[vlab, below right] at (C) {$C$};
  \node[vlab] at ($(O)+(167:0.42)$) {$O$};
\end{tikzpicture}
\figcap{13--11}\end{bkfigure}}

% --------------------------------------------------------------- Fig 13-12
% Thm 13-5, CIRCUMCENTER.  NOTE: this figure is the circumcenter, not a second
% incenter -- the dashed lines leave the MIDPOINTS of the sides, not the
% vertices, and O falls outside the triangle because angle B is obtuse.
% O is built as an intersection of two perpendicular bisectors.
\newcommand{\FIGXIIITWELVE}{\begin{bkfigure}[\linewidth]
\begin{tikzpicture}[scale=1.300]
  \coordinate (A) at (0,0);
  \coordinate (B) at (3.20,0);
  \coordinate (C) at (3.75,2.75);
  \coordinate (Mab) at ($(A)!0.5!(B)$);
  \coordinate (Mbc) at ($(B)!0.5!(C)$);
  \coordinate (Mac) at ($(A)!0.5!(C)$);
  % perpendicular bisector of AB: through Mab, perpendicular to AB
  \coordinate (Pab) at ($(Mab)+(0,1)$);
  % perpendicular bisector of BC: through Mbc, perpendicular to BC
  \coordinate (Pbc) at ($(Mbc)!1!90:(C)$);
  \coordinate (O) at (intersection of Mab--Pab and Mbc--Pbc);
  \draw[given] (A) -- (B) -- (C) -- cycle;
  % Re-drawn 2026-08-17.  The book draws each perpendicular bisector SHORT: it
  % starts at the midpoint of its side and stops in a small stub just past O.
  % The old drawing ran all three clean through the triangle and out the far
  % side into empty space, which is what Caleb flagged.
  \draw[auxlong] (Mab) -- ($(Mab)!1.26!(O)$);
  \draw[auxlong] (Mbc) -- ($(Mbc)!1.22!(O)$);
  \draw[auxlong] (Mac) -- ($(Mac)!2.25!(O)$);
  % the cross-strokes at the midpoints -- these are what identify the three
  % lines as PERPENDICULAR bisectors rather than as arbitrary cevians
  \draw[tick] ($(Mab)!0.075!90:(B)$) -- ($(Mab)!-0.075!90:(B)$);
  \draw[tick] ($(Mbc)!0.0856!90:(C)$) -- ($(Mbc)!-0.0856!90:(C)$);
  \XIIIstar{O}
  \node[vlab, below left] at (A) {$A$};
  \node[vlab, below right] at (B) {$B$};
  \node[vlab, above right] at (C) {$C$};
  \node[vlab] at ($(O)+(40:0.40)$) {$O$};
\end{tikzpicture}
\figcap{13--12}\end{bkfigure}}

% --------------------------------------------------------------- Fig 13-13
% Ex 13-4 #1.  CD is PARALLEL to the base: C and D are taken at the SAME ratio
% down the two sides, which forces CD || AB exactly.  P is the midpoint of CD.
\newcommand{\FIGXIIITHIRTEEN}{\begin{bkfigure}[\linewidth]
\begin{tikzpicture}[scale=1.300]
  \coordinate (A) at (0,0);
  \coordinate (Bs) at (3.64,-0.18);
  \coordinate (T) at (0.73,2.09);
  \coordinate (C) at ($(T)!0.4785!(A)$);
  \coordinate (D) at ($(T)!0.4785!(Bs)$);
  \coordinate (P) at ($(C)!0.5!(D)$);
  \draw[given] (A) -- (Bs) -- (T) -- cycle;
  \draw[given] (C) -- (D);
  \fill (P) circle (0.055);
  \node[vlab, left] at (C) {$C$};
  \node[vlab, above right] at (D) {$D$};
  \node[vlab, below] at (P) {$P$};
\end{tikzpicture}
\figcap{13--13}\end{bkfigure}}

% --------------------------------------------------------------- Fig 13-14
% Ex 13-4 #4.  O is the intersection of the perpendicular bisector of AB with
% the perpendicular to BC at B -- so OB _|_ BC exactly, and |OA| = |OB|.
% Built the other way round for stability: put O at the centre, take A and B
% on the circle, then erect BC perpendicular to OB at B.
\newcommand{\FIGXIIIFOURTEEN}{\begin{bkfigure}[\linewidth]
\begin{tikzpicture}[scale=1.300]
  \def\rr{1.50}
  \coordinate (O) at (0,0);
  \coordinate (A) at (203:\rr);
  \coordinate (B) at (-19:\rr);
  \coordinate (Q) at (59:\rr);
  % C on the perpendicular to OB at B, on the far side from the centre
  \coordinate (C) at ($(B)!1.05!90:(O)$);
  \draw[given] (O) circle (\rr);
  \draw[given] (A) -- (B);
  \draw[given] (B) -- (C);
  % perpendicular bisector of AB (through O), and the radius OB
  \coordinate (Mab) at ($(A)!0.5!(B)$);
  % drawn as in the print: a short overshoot above O, then down past the
  % chord AB to the foot of the circle
  \draw[aux] ($(O)!-0.89!(Mab)$) -- ($(O)!2.85!(Mab)$);
  \draw[aux] (B) -- ($(O)!-0.62!(B)$);
  \XIIIplus{O}
  \fill (Q) circle (0.055);
  \node[vlab, left] at (A) {$A$};
  \node[vlab, right] at (B) {$B$};
  \node[vlab] at ($(O)+(216:0.34)$) {$O$};
  \node[vlab, above right] at (Q) {$Q$};
  \node[vlab, below] at (C) {$C$};
\end{tikzpicture}
\figcap{13--14}\end{bkfigure}}

% ------------------------------------------------------- Figs 13-15, 13-16
% Printed side by side on p.233.
%
% 13-15, Thm 13-6 CENTROID: D, E are midpoints of AB, BC; O is the crossing of
% the medians CD and AE; G, F are the midpoints of OA, OC (so that D, E, F, G
% are the parallelogram vertices of the proof outline).
%
% 13-16, Thm 13-7 ORTHOCENTER: through each vertex, the parallel to the
% opposite side.  That makes A, B, C exactly the midpoints of the sides of
% A'B'C', so A' = B+C-A, B' = C+A-B, C' = A+B-C -- computed, not eyeballed.
% O is the orthocentre of ABC (= circumcentre of A'B'C').
\newcommand{\FIGXIIIFIFTEENSIXTEEN}{\begin{bkfigure}[\linewidth]
\hfill
\begin{minipage}{0.43\linewidth}\centering
\begin{tikzpicture}[scale=0.92]
  \coordinate (A) at (0,0);
  \coordinate (B) at (2.00,2.14);
  \coordinate (C) at (4.10,-0.10);
  \coordinate (D) at ($(A)!0.5!(B)$);
  \coordinate (E) at ($(B)!0.5!(C)$);
  \coordinate (O) at (intersection of A--E and C--D);
  \coordinate (G) at ($(A)!0.5!(O)$);
  \coordinate (F) at ($(O)!0.5!(C)$);
  \draw[given] (A) -- (B) -- (C) -- cycle;
  \draw[given] (A) -- (E);
  \draw[given] (C) -- (D);
  \draw[given] (D) -- (E);
  \fill (G) circle (0.055);
  \fill (F) circle (0.055);
  \node[vlab, left] at (A) {$A$};
  \node[vlab, above] at (B) {$B$};
  \node[vlab, right] at (C) {$C$};
  \node[vlab, above left] at (D) {$D$};
  \node[vlab, above right] at (E) {$E$};
  \node[vlab] at (1.38446,0.20798) {$G$};
  \node[vlab] at (2.75246,0.16316) {$F$};
  \node[vlab] at ($(O)+(269:0.32)$) {$O$};
\end{tikzpicture}
\figcap{13--15}
\end{minipage}
\hfill
\begin{minipage}{0.55\linewidth}\centering
\begin{tikzpicture}[scale=0.82]
  \coordinate (A) at (0,0);
  \coordinate (B) at (1.10,0.95);
  \coordinate (C) at (2.86,-0.05);
  \coordinate (Ap) at ($(B)+(C)-(A)$);
  \coordinate (Bp) at ($(C)+(A)-(B)$);
  \coordinate (Cp) at ($(A)+(B)-(C)$);
  % orthocentre: intersection of two altitudes, each a genuine perpendicular
  \coordinate (Fa) at ($(B)!(A)!(C)$);
  \coordinate (Fb) at ($(A)!(B)!(C)$);
  \coordinate (O) at (intersection of A--Fa and B--Fb);
  % Re-drawn 2026-08-17 against the print.  The arrowheads had been put on the
  % WRONG lines: the book arrows the three SOLID sides (two tips converging on
  % C, one on A) and leaves the dashed altitudes AO and CO plain, with a single
  % arrowhead at O belonging to the altitude through B.  Ours arrowed AO and CO
  % into A and C, so three or four heads piled onto each vertex and blotted it
  % -- which is also what made the straight line C'-A-B' look kinked at A.
  \draw[auxlong,bkarrows] ($(Cp)!-0.04!(Ap)$) -- ($(Ap)!-0.04!(Cp)$);
  \draw[auxlong] (Cp) -- (Bp);
  \draw[auxlong] (Ap) -- (Bp);
  \draw[given,bkarrow] (B) -- (A);
  \draw[given,bkarrow] (B) -- (C);
  \draw[given,bkarrow] (A) -- (C);
  \draw[auxlong] (A) -- (O);
  \draw[auxlong] (C) -- (O);
  \draw[auxlong,bkarrow] (Fb) -- (O);
  % the print ticks AC where the altitude through B meets it
  \draw[tick] ($(Fb)!0.055!90:(C)$) -- ($(Fb)!-0.055!90:(C)$);
  \node[vlab, below left] at (A) {$A$};
  \node[vlab] at ($(B)+(44:0.46)$) {$B$};
  \node[vlab, below right] at (C) {$C$};
  \node[vlab, right] at ($(Ap)!-0.04!(Cp)$) {$A\XIIIpr$};
  \node[vlab, below] at (Bp) {$B\XIIIpr$};
  \node[vlab, left] at ($(Cp)!-0.04!(Ap)$) {$C\XIIIpr$};
  \node[vlab, above] at (O) {$O$};
\end{tikzpicture}
\figcap{13--16}
\end{minipage}
\hfill
\end{bkfigure}}

% --------------------------------------------------------------- Fig 13-17
% Ex 13-5 *3: the data for an SSA construction -- the angle at A, with the two
% given lengths tick-marked.  Reproduced as the book draws it: the second side
% of the angle at C runs free (the third vertex is not yet located), and the
% triple-ticked segment is drawn across from that free side to CA.
% Re-drawn 2026-08-17.  Caleb settled this figure: "YOU DID THIS CORRECT BESIDES
% THE 1 TICK NOT BEING PERPENDICULAR TO THE LINE CA."  The single tick had been
% struck at 115 degrees while CA runs at -43.6, i.e. only 25 degrees off lying
% ALONG the line -- which is what smeared it into the triple tick below it.
% Every tick group is now struck perpendicular to the segment it measures, the
% three ticks are shortened and spread so they stay countable, and the whole
% figure is re-proportioned to the print: angle at C = 26 degrees, the free ray
% 1.36|CA| long, |AB| = 1.745|CA|, the polyline C-A-B turning 39 degrees at A,
% the single tick at 0.468 of CA and the transversal foot at 0.70 (it used to be
% at 0.85, which crowded the two tick groups together).
\newcommand{\FIGXIIISEVENTEEN}{\begin{bkfigure}[\linewidth]
\begin{tikzpicture}[scale=1.450]
  \coordinate (C) at (0,1.30);
  \coordinate (A) at ($(C)+(-43.6:1.55)$);
  \coordinate (Cf) at ($(C)+(-17.6:2.108)$);   % free end of the second ray
  \coordinate (B) at ($(A)+(-4.9:2.705)$);
  \draw[given] (C) -- (A);
  \draw[given] (A) -- (B);
  \draw[given] (C) -- (Cf);
  % the triple-ticked transversal, from a point on CA across to the free ray
  \coordinate (Sf) at ($(C)!0.70!(A)$);
  \coordinate (S)  at ($(C)!0.51!(Cf)$);
  \draw[given] (Sf) -- (S);
  % ONE tick on CA, PERPENDICULAR to CA
  \coordinate (T1) at ($(C)!0.468!(A)$);
  \draw[tick] ($(T1)!0.1394!90:(A)$) -- ($(T1)!-0.1394!90:(A)$);
  % THREE ticks on the transversal, perpendicular to it, grouped toward the CA
  % end as the print groups them
  \coordinate (Qt) at ($(Sf)!0.144!90:(S)$);
  \foreach \t in {0.28,0.41,0.54}
    \draw[tick] ($($(Sf)!\t!(S)$)+(Qt)-(Sf)$) -- ($($(Sf)!\t!(S)$)-(Qt)+(Sf)$);
  % TWO ticks on AB, perpendicular to AB
  \coordinate (Qb) at ($(A)!0.0388!90:(B)$);
  \foreach \t in {0.47,0.53}
    \draw[tick] ($($(A)!\t!(B)$)+(Qb)-(A)$) -- ($($(A)!\t!(B)$)-(Qb)+(A)$);
  \node[vlab, above left] at (C) {$C$};
  \node[vlab, below left] at (A) {$A$};
  \node[vlab, below right] at (B) {$B$};
\end{tikzpicture}
\figcap{13--17}\end{bkfigure}}


% ---- local macros (hoist into book preamble) ----
% The "*" on _locus_ on p.227 is a GLOSSARY footnote, not the starred-exercise
% convention, so \starnote is wrong here.  As with ch02's \IIstarnote and
% ch09's \IXstarnote, \footnotetext alone would print the footnote counter;
% the book marks this note with the "*" itself and no number.
\newcommand{\XIIIlocusnote}{\begingroup\renewcommand{\thefootnote}{}%
  \footnotetext{*\emph{Locus} is a Latin word meaning, literally, \emph{place}
  in the sense that we use the word \emph{location}.}%
  \endgroup}
% ---- end local macros ----

\begin{document}
\setcounter{chapter}{13}

\chapter*{\raggedright\Huge LOCI AND SETS}
\markboth{LOCI AND SETS}{}
\addcontentsline{toc}{chapter}{13.\ Loci and Sets}
\vspace{-1em}\noindent\rule{\linewidth}{0.6pt}\vspace{1em}
\figurekey

%========================================================= 13-1
\section*{13--1\quad TERMINOLOGY}
\addcontentsline{toc}{section}{13--1\quad Terminology}

In this chapter we use the term \emph{locus}*\XIIIlocusnote{} (plural,
\emph{loci}), which we inherit from the days when Latin was the universal
language of science. (Even in the first part of the 19th century, most
mathematics was published in Latin.) As will be clear from the definition given
below, a locus is just a set of points. The language may be met in later work,
so we devote this short chapter to it. We also collect a number of geometric
constructions in this chapter.

\begin{definition}
A locus of points is the set of all points, and only those points, which
satisfy some given condition.
\end{definition}

We may now use the language, ``The locus of points $P$ such that\ldots.'' This
is just another way of saying, ``The set of all points $P$ such that\ldots.''

Examples:

(a) The \emph{locus} of points at a given distance $r$ from a given point is
obviously a circle.

(b) The \emph{locus} of points at a given distance $d$ from a given line $l$ is
a pair of lines parallel to $l$, one on each side of $l$ (Fig.~13--1).

\FIGXIIIONE

Example (b) requires a bit of proof, which is omitted. We will discuss such
proofs in the next section. This locus could also be described as the \emph{set}
of all points at distance $d$ from $l$.

\FIGXIIITWO

(c) The \emph{locus} of points equidistant from two parallel lines is a line
between the given lines.

An equivalent way of describing (c) would be to say: the \emph{set} of all
points equidistant from two parallel lines (Fig.~13--2). Again, the fact that
locus (c) is a line requires proof.

\exercisegroup{13--1}

In the following exercises, state what you believe the locus to be. Also
describe this locus as a \emph{set}.

\begin{exlist}
\item What is the locus of points $P$ such that $\sg{PP_0} = 4$, where $P_0$ is
      a given point?
\item What is the locus of all points at a distance 1 from a line segment three
      units long (Fig.~13--3)? (By the distance from a point to a line segment,
      we mean the shortest distance from the point to any point of the
      segment.)
\exfig{\FIGXIIITHREE}
\item What is the locus of all points at a distance four units from the sides of
      a unit square? a square 12 units on a side? eight units?
\item A circle of radius five units is given. What is the locus of all points at
      a distance two units from the given circle?
\item A triangle has a fixed base four units long (Fig.~13--4). What is the
      locus of all possible positions of the vertex $V$ if the area is four
      units? if the area is $A$ units?
\exfig{\FIGXIIIFOUR}
\end{exlist}

%========================================================= 13-2
\section*{13--2\quad SOME LOCUS THEOREMS}
\addcontentsline{toc}{section}{13--2\quad Some Locus Theorems}

In many locus problems the set of points can be characterized in more than one
way, or we can say that two locus statements may describe the same set of
points. When this is the case, there is the following to be shown: (1)~Each
point of the first locus is also a point of the second locus. (The first set is
a \emph{subset} of the second set.) (2)~Each point of the second locus is also a
point of the first locus. (The second set is a subset of the first set.) If each
is a subset of the other, the two sets are the same.

\FIGXIIIFIVE

\begin{theorem}
The locus of points $P$ equidistant from two given points, $A$ and $B$, is the
set of points on the perpendicular bisector of the segment $AB$ (Fig.~13--5).
\end{theorem}

To make the proof we must show two things: (1)~that if a point $P$ is
equidistant from $A$ and $B$, then it is on the perpendicular bisector, and
(2)~that if a point is on the perpendicular bisector, then it is equidistant
from $A$ and $B$. These should be easy for you to prove now; in fact, both (1)
and (2) were proven previously as theorems or exercises.

A brief proof will be sketched here, although your teacher may require you to
write out the details. To prove (1), let $C$ denote the midpoint of $AB$. Then
if $P \neq C$, $\tri{PAC} \cong \tri{PBC}$. (Why?) Thus $\ang{PCA} \cong
\ang{PCB}$, and $P$ is on the perpendicular bisector. (How do we argue if
$P = C$?) To prove (2), suppose that $P$ is on the perpendicular bisector and
$P \neq C$. Then $\tri{PCA} \cong \tri{PCB}$ (why?) and $PA \cong PB$. (How do
we argue if $P = C$?)

A similar theorem concerns angle bisectors.

\begin{theorem}
The locus of points interior to an angle and equidistant from the sides of the
angle is the set of points on the bisector of the angle, excluding the vertex
(Fig.~13--6).
\end{theorem}

\FIGXIIISIX

\noindent\emph{Proof sketch.} Suppose that $P$ is interior to the angle and
equidistant from the sides. Let $A$ and $B$ be the feet of perpendiculars from
$P$ to the sides of the angle. (How do you know that perpendiculars from $P$
meet the sides of the angle?) Then $\tri{POA}$ and $\tri{POB}$ are right
triangles with two congruent sides and are congruent. Hence $\ang{POA} \cong
\ang{POB}$, and $OP$ is the angle bisector. Conversely, if $P$ is on the
bisector, then $P \neq O$, $PA$ and $PB$ are perpendicular to the sides of the
angle, and the right triangles $POA$ and $POB$ are congruent. Hence
$\sg{PA} = \sg{PB}$.

\begin{theorem}
Let $AB$ be a given segment. The set of all points $P$ such that $\tri{APB}$ is
right, with $AB$ the hypotenuse, is the circle on $AB$ as diameter, excluding
points $A$ and $B$ (Fig.~13--7).
\end{theorem}

\FIGXIIISEVEN

Let $O$ be the midpoint of $AB$. It is necessary to show that if $P$ is a point
such that $AP \perp PB$, then $P$ is on the circle whose center is $O$ and which
contains $A$ and $B$.

Conversely, if $P$ is on the circle, and different from $A$ and $B$, it is
necessary to show that $AP \perp PB$. But this follows from theorems on the
measure of inscribed angles. (There are other proofs, too.) The details are left
to the student.

\exercisegroup{13--2}

\begin{exlist}
\item What is the set of all points that are centers of circles tangent to both
      sides of an angle?
\item What is the locus of all points exterior to a given circle and at a
      distance $d$ from the circle?
\item A barn has dimensions $30 \times 50$ ft. A cow is tethered to a corner of
      the barn by a rope 60 ft long. What is the locus of points where the cow
      may graze?
\item[\stex 4.] In Fig.~13--8, points $A$ and $B$ are fixed. What is the locus
      of all points $P$, on a given side of the line $AB$, such that
      $\ang{APB}$ is $50^\circ$?
\exfig{\FIGXIIIEIGHT}
\item[5.] What is the locus of points equidistant from three noncollinear
      points?
\item[6.] As in Fig.~13--9, sketch the path of a point on the rim of a wheel
      rolling on a line. The curve obtained is called a \emph{cycloid}.
\exfig{\FIGXIIININE}
\item[\stex 7.] With the same wheel as in Exercise 6, sketch the locus if $P$ is
      on a spoke of the wheel; on an extension of a spoke.
\item[8.] A small circle of radius $R/4$ rolls on the inside of a circle of
      radius $R$ (Fig.~13--10). What is the path of a point on the small circle?
      (The curve is called a \emph{hypocycloid}.)
\exfig{\FIGXIIITEN}
\item[\stex 9.] What are the various possibilities that arise from Exercise 8 if
      the radius of the small circle is $xR$? Investigate the locus if
      $x = \frac{1}{2}$; $x = \frac{1}{3}$; $x$ is a rational number; $x$ is an
      irrational number.
\end{exlist}

%========================================================= 13-3
\section*{13--3\quad UNIONS AND INTERSECTIONS}
\addcontentsline{toc}{section}{13--3\quad Unions and Intersections}

Some loci are just unions or intersections of other simpler loci. To be precise,
we make a definition.

\begin{definition}
Given two sets of points $S_1$ and $S_2$, the union of the two sets is the set of
all points of $S_1$ and $S_2$. The \emph{intersection} of $S_1$ and $S_2$ is the
set of points that are common to $S_1$ and $S_2$.
\end{definition}

\exercisegroup{13--3}

\begin{multicols}{2}
\begin{exlist}
\item Describe an angle as a union of two sets.
\item Describe an angle and its interior as the intersection of two sets.
\item Describe the interior of a triangle as the intersection of three sets; as
      the intersection of two sets.
\item Describe a triangle as the union of three sets.
\item Describe a quadrilateral as the union of four sets. Describe the interior
      of a simple convex quadrilateral as the intersection of two sets; as the
      intersection of four sets.
\end{exlist}
\end{multicols}

\begin{theorem}
The bisectors of the angles of a triangle intersect in one point (called the
\emph{incenter}), which is equidistant from the three sides (Fig.~13--11).
\end{theorem}

\FIGXIIIELEVEN

\noindent\emph{Proof sketch.} The bisectors of angles $A$ and $B$ must intersect
in one point $O$. (Why?) Then this point $O$ is equidistant from $AB$ and from
$BC$. (Why?) Likewise, $O$ is equidistant from $AB$ and $AC$. Therefore, $O$ is
on the bisector of $\ang{C}$. (Why?) This proves the theorem.

The point $O$ is called the \emph{incenter} because it is the center of a circle
\emph{inscribed} in the triangle. One may also state the theorem by saying that
\emph{the bisectors of the angles are concurrent}.

\FIGXIIITWELVE

\begin{theorem}
The perpendicular bisectors of the sides of a triangle intersect in one point
(called the \emph{circumcenter}), which is equidistant from the three vertices
(Fig.~13--12).
\end{theorem}

The proof is left to the student.

\exercisegroup{13--4}

\begin{exlist}
\item Consider all lines that are parallel to the base of a triangle and that
      intersect the other sides in points $C$ and $D$ (Fig.~13--13). Show that
      the locus of midpoints of such segments $CD$ is the median to the base
      (except for one point!).
\exfig{\FIGXIIITHIRTEEN}
\item Show how to inscribe a circle in a triangle. State why your construction
      is correct. Why is there a unique inscribed circle?
\item Show how to circumscribe a circle about a triangle. State why there is a
      unique circumscribing circle.
\item In Fig.~13--14, $\ang{ABC}$, neither right nor straight, is given. The
      point $O$ is the intersection of the perpendicular bisector of $AB$ and
      the line perpendicular to $BC$ through $B$. Prove that there is a circle
      with center $O$ containing $A$ and $B$. If $Q$ is a point of this circle
      and is on the same side of $AB$ as $O$, prove that
      $\ang{AQB} \cong \ang{ABC}$.
\exfig{\FIGXIIIFOURTEEN}
\end{exlist}

\begin{theorem}
The medians of a triangle intersect in one point (called the \emph{centroid}),
which is two thirds of the distance from each vertex to the midpoint of the
opposite side.
\end{theorem}

\noindent\emph{Outline of proof.} (Details are left as an exercise.) (1)~Any two
medians, say $CD$ and $AE$ in Fig.~13--15, intersect. Call this intersection
point $O$. (2)~Choose $F$ and $G$ as the midpoints of $OC$ and $OA$,
respectively. (3)~Show that $D$, $E$, $F$, and $G$ are vertices of a
parallelogram. (4)~Thus $AG \cong GO \cong OE$, and $FC \cong FO \cong OD$.
(5)~The point $O$ has, therefore, the required distance properties. Since there
is but one such point on each median, the medians are concurrent.

The centroid is the ``balance point'' in the following sense. A sheet of metal
of uniform thickness, the shape of the triangle, would balance if supported at
$O$. Exercise 1 of Group 13--4 may help you to see why this is so.

\FIGXIIIFIFTEENSIXTEEN

\begin{theorem}
The altitudes of a triangle intersect in one point (called the
\emph{orthocenter}).
\end{theorem}

\noindent\emph{Outline of proof.} (Details are left as an exercise.)
(1)~Through each vertex, consider the line that is parallel to the opposite side
of the given $\tri{ABC}$. These lines will intersect to form $\tri{A'B'C'}$, as
in Fig.~13--16. (2)~Show that $AC'BC$, $ABCB'$, and $ABA'C$ are parallelograms,
and hence observe that $C'B \cong BA'$, $C'A \cong AB'$, and $B'C \cong CA'$.
(3)~Show that the circumcenter $O$ of $\tri{A'B'C'}$ is the point of
intersection of the altitudes of $\tri{ABC}$.

\exercisegroup{13--5}

\begin{exlist}
\item Construct a triangle, given the lengths of its base, altitude, and one
      side.
\item Construct a triangle, given the base, the length of the altitude to this
      base, and one of the base angles.
\item[\stex 3.] As in Fig.~13--17, construct a triangle, given the lengths of
      two sides and an angle other than the included angle. [\emph{Hint:} Use
      Exercise 4 in Group 13--4.]
\exfig{\FIGXIIISEVENTEEN}
\item[4.] Construct a triangle, given the lengths of an altitude to one side and
      the other two sides.
\item[5.] Construct a triangle, given the midpoints of the three sides.
\item[6.] Construct a triangle, given the length of the altitude to the base and
      the two base angles.
\item[7.] Given a circle and a line segment $AB$, construct the locus of all
      points $P$ such that tangents of length $\sg{AB}$ can be drawn to the
      circle from $P$.
\item[8.] Construct a right triangle with a given segment as hypotenuse and
      having a side of given length.
\item[9.] Construct a circle tangent to a given line at a given point and having
      a given radius.
\item[10.] Construct the locus of all points $P$ such that there is a circle,
      with given radius and center at $P$, which is tangent to a given circle.
\item[11.] Construct all circles with a given radius which are tangent to two
      given circles of unequal radii.
\item[12.] Construct a line tangent to a given circle and perpendicular to a
      given line.
\item[13.] Given a circle and a line segment $AB$, construct the set of all
      points $M$ such that $M$ is the midpoint of a chord (of the circle)
      congruent to $AB$.
\item[14.] Given an arc of a circle, construct the center of the circle.
\item[15.] Construct an isosceles triangle, given the base, and the radius of
      the circumscribed circle.
\item[16.] Construct an isosceles triangle, given the base, and the radius of
      the inscribed circle.
\item[17.] Construct an isosceles triangle, given the base and the vertex angle.
\item[18.] Construct a triangle, given two sides and the length of the median to
      one of these sides.
\item[\stex 19.] Let $P$ be a given point not on a given circle whose center is
      $O$. Consider all segments $PX$ such that $X$ is on the circle. Show that
      the locus of midpoints $M$ of these segments is a circle whose diameter is
      a segment on the line through $P$ and $O$. Construct this circle.
\item[\stex 20.] Given three segments, construct a triangle two of whose sides
      are congruent to two of these segments, and such that the median to the
      other side is congruent to the third segment.
\end{exlist}

%========================================================= Review
\section*{REVIEW OF CHAPTER 13}
\addcontentsline{toc}{section}{Review of Chapter 13}
\markboth{REVIEW OF CHAPTER 13}{}

In earlier chapters we were using the phrases \emph{set of points} and
\emph{geometric figure} interchangeably. In this chapter we introduced a third
term, \emph{locus}, which has exactly the same meaning as the earlier
expressions.

It is important that you understand what is involved in establishing that two
different ``rules'' determine the same locus. A statement like

\begin{quote}
\emph{the set of points having property $p$ is the same set as that having
property $q$}
\end{quote}

\noindent must be established by proving two things:

(1) Every point in the first set is also in the second. That is, every point
having property $p$ also has property $q$.

(2) Every point in the second set is also in the first. That is, every point
having property $q$ also has property $p$.

% The book heads this "Review Exercises", with no "Exercise"/"Exercise Group"
% prefix, so \exercise{} (which prepends "Exercise") is wrong here.
\par\medskip\begin{center}\bfseries Review Exercises\end{center}\par\nopagebreak

\begin{exlist}
\item If $AB$ is a fixed segment, are the properties
      \begin{quote}
      $p$: $\ang{AXB}$ is acute
      \end{quote}
      \noindent and
      \begin{quote}
      $q$: $X$ is in the exterior of the circle on $AB$ as diameter
      \end{quote}
      \noindent equivalent to each other?
\item Restate Exercise 1 in conventional locus terminology.
\item If $R$ is the set of points $X$ such that, for a fixed segment $AB$,
      % As printed: the book sets the right-hand side of this first relation
      % WITHOUT an overbar, though the two relations below it carry one.
      % Verified at 320 dpi (PDF 249); "AB," opens the line, so nothing is
      % clipped.  Reproduced as printed -- see ch13-UNCERTAIN.md.
      $\sg{AX} + \sg{XB} = AB$, and $S$ is the set of points between $A$
      and $B$, prove that $R = S$. How else can you describe the locus of points
      $X$ such that $\sg{AX} - \sg{XB} = \sg{AB}$? such that
      $\sg{BX} - \sg{XA} = \sg{AB}$?
\item Let $AB$ and $CD$ be two given segments. Describe in a second way the set
      of points $X$ such that
      \begin{quote}
      $\ang{AXB}$ is obtuse and $\ang{CXD}$ is right;\\
      $\ang{AXB}$ is obtuse or $\ang{CXD}$ is right;\\
      $\ang{AXB}$ is acute and $\ang{CXD}$ is obtuse;\\
      $\ang{AXB}$ is acute or $\ang{CXD}$ is obtuse.
      \end{quote}
\end{exlist}

\end{document}
